When an agent issues several tool calls at once, or a team of large language model (LLM) agents splits a task, each decision must take its own role without seeing the others. We study where such roles come from, with diffusion LLMs (dLLMs) as a test bed: a dLLM agent fills the value slots of a fixed call skeleton, and its decoding schedule sets which slots decide together and what each sees. We compare these slots with teams of 7B LLM agents that each make one call. On the parallel calls of the Berkeley Function Calling Leaderboard, slots that decide in the same step take the entities in the order the request mentions them: committing 16 tokens per step instead of one costs the dLLM Dream 5.3 points of set accuracy. Simultaneous agents that see only their own call rarely recover this order, whether numbered (28.8\%, against 81.7\% for the slots) or told the rule (48.1\%), even where the whole skeleton tells a slot no more than its number; agents that take turns come within 4 points of one agent making all calls. When nothing orders the choices, simultaneous agents collide, and sampling does not prevent it. Slot length is a strong cue: wrong lengths corrupt values and bind them to the wrong calls even with one token per step, partly through two interface choices we isolate, a pre-written closing delimiter and a type mask. These errors look like failures of parallel decoding, so evaluations of skeleton-based dLLM agents should control slot lengths before blaming parallelism.
